\section{Proper-model norms of arbitrary original residue-fibre vectors}

\begin{definition}[Transport of an original residue-fibre vector]
\label{def:original-residue-vector-transport}
\lean{EnvelopeProperty.Berkovich.SchemeFramedSectionPullback.pullbackResidueSection_eq_affinePullbackUnitTop, EnvelopeProperty.Berkovich.SchemeFramedSectionPullback.transportedResidueSection}
\uses{berkovich:actual-residue-module-pullback-iso,def:fc-induced-native-residue-frame}
\leanok
For any scheme morphism $g:Y\to Z$, a point $y\in Y$, an original
module $M$ on $Z$, and an actual isomorphism $\alpha:g^*M\simeq N$,
transport any vector $s$ in the residue fibre $M(g(y))$ first along
the actual field homomorphism $\kappa(g(y))\to\kappa(y)$ and then
along the residue pullback of $\alpha$. Denote the resulting vector
of $N(y)$ by $s_Y$. No extension of $s$ to a model section is required.
\end{definition}

\begin{lemma}[The induced frame preserves every residue coefficient]
\label{lem:original-residue-vector-frame-coefficient}
\lean{EnvelopeProperty.Berkovich.SchemeFramedSectionPullback.transportedResidueSection_frameSection, EnvelopeProperty.Berkovich.SchemeFramedSectionPullback.transportedResidueSection_coefficient, EnvelopeProperty.Berkovich.SchemeFramedSectionPullback.transportedResidueSection_valuation}
\uses{def:original-residue-vector-transport,def:fc-induced-native-residue-frame,lem:fc-native-unit-section-morphism}
\leanok
Let $e$ be any actual residue-fibre frame of $M$ at $g(y)$, and let
$e_Y$ be its canonical induced frame on $N(y)$. The coefficient
$e_Y(s_Y)$ is the image of the original coefficient $e(s)$ under
the same actual residue-field homomorphism. Consequently its value
under the valuation at $y$ equals the value of $e(s)$ at the image
valued point.
\end{lemma}
\begin{proof}
The entire composition of the residue-module pullback comparison,
the residue pullback of $\alpha$, and the induced frame is the
pullback of $e$ followed by the canonical structure-sheaf comparison.
Apply this identity to the actual residue-field pullback of $s$.
Naturality of that pullback for $e$ identifies the resulting vector
with the pullback of $e(s)$. The structure-sheaf comparison maps
this vector to its image under the actual map of spectra. Naturality
of global sections of an affine spectrum identifies this image with
the stated field image. The valuation of the image valued point is
defined by composition with this same field map, giving the equality
of values, including when the vector is zero.
\end{proof}

\begin{theorem}[The proper-model norm of every original residue vector]
\label{thm:proper-model-arbitrary-residue-vector-norm}
\lean{EnvelopeProperty.Berkovich.RelativeSchemeValuedPoint.modelLineNorm_framedResiduePullback, EnvelopeProperty.Berkovich.RelativeSchemeValuedPoint.exists_sourceResidueFrame_modelLineNorm_residueSections}
\uses{lem:original-residue-vector-frame-coefficient,prop:generated-line-generic-frame-norms}
\leanok
Assume $Z\to\operatorname{Spec}R$ is proper, the base valuation is
bounded by one, and $M$ is invertible. For any valued point of an
arbitrary source $Y\to Z$, every vector $s\in M(g(y))$ satisfies
\[
 \|s\|_M=|e_Y(s_Y)|\,\|e^{-1}(1)\|_M.
\]
In particular one induced source residue frame computes the norm
of every vector simultaneously with no extra factor. The source
scheme need not be proper and its valuation need not be bounded
on its whole coordinate ring.
\end{theorem}
\begin{proof}
In the one-dimensional target residue fibre the actual frame gives
$s=e(s)e^{-1}(1)$. Homogeneity of the original proper-model norm
therefore gives the product of the value of $e(s)$ and the norm of
$e^{-1}(1)$. Substitute the coefficient-value identity of the preceding
lemma to obtain the displayed formula. Properness provides the actual
valuation-ring center of the image point, and invertibility gives a
frame over its local valuation ring. The residue pullback of that
frame has inverse-unit section of norm one, because its coefficient
in this same valuative frame is one. Use this target frame and its
canonical induced source frame. The formula then holds with unit
factor for every residue vector at once.
\end{proof}
